\chapter{N-Deklination}
\label{cha:NDeklination}

Die N-Deklination betrifft eine spezifische Gruppe maskuliner Nomen, die in allen Fällen außer Nominativ Singular eine Endung erhalten.

\section{Nomen der N-Deklination}
\label{sec:NDeklinationNomen}

\subsection{Nomen auf -e}
\begin{itemize}
    \item der Junge $\rightarrow$ dem Jungen, den Jungen
    \item der Kunde $\rightarrow$ dem Kunden, den Kunden
    \item der Kollege $\rightarrow$ dem Kollegen, den Kollegen
    \item der Arbeiter $\rightarrow$ dem Arbeiter, die Arbeiter
\end{itemize}

\subsection{Nomen auf -ant, -ent, -ist}
\begin{itemize}
    \item der Praktikant $\rightarrow$ dem Praktikanten
    \item der Student $\rightarrow$ dem Studenten
    \item der Polizist $\rightarrow$ dem Polizisten
    \item der Journalist $\rightarrow$ dem Journalisten
\end{itemize}

\subsection{Ausnahmen (unregelmäßige N-Deklination)}
\begin{itemize}
    \item der Herr $\rightarrow$ dem Herrn, die Herren
    \item der Mensch $\rightarrow$ dem Menschen, die Menschen
    \item der Nachbar $\rightarrow$ dem Nachbarn, die Nachbarn
    \item der Bauer $\rightarrow$ dem Bauern, die Bauern
    \item der Held $\rightarrow$ dem Helden, die Helden
    \item der Prinz $\rightarrow$ dem Prinzen, die Prinzen
    \item der Kunde $\rightarrow$ dem Kunden, die Kunden
    \item der Experte $\rightarrow$ dem Experten, die Experten
    \item der Architekt $\rightarrow$ dem Architekten, die Architekten
\end{itemize}

\section{Deklinationstabelle}
\label{sec:NDeklinationTabelle}

\begin{center}
\begin{tabular}{|l|c|c|}
\hline
\textbf{Kasus} & \textbf{Singular} & \textbf{Plural} \\ \hline
Nominativ & der Architekt & die Architekten \\ \hline
Akkusativ & den Architekten & die Architekten \\ \hline
Dativ & dem Architekten & den Architekten \\ \hline
Genitiv & des Architekten & der Architekten \\ \hline
\end{tabular}
\end{center}

\section{Übungen zur N-Deklination}
\label{sec:NDeklinationUebungen}

\begin{tcolorbox}[title={Aufgabe}]
Deklinieren Sie die folgenden Nomen im Dativ Singular und Plural.
\end{tcolorbox}

\begin{enumerate}[label=\textbf{Aufgabe \arabic*}.]
    \item der Architekt (Dativ Sg.: \underline{\hspace{1cm}}, Pl.: \underline{\hspace{1cm}})
    \item der Kunde (Dativ Sg.: \underline{\hspace{1cm}}, Pl.: \underline{\hspace{1cm}})
    \item der Student (Dativ Sg.: \underline{\hspace{1cm}}, Pl.: \underline{\hspace{1cm}})
    \item der Herr (Dativ Sg.: \underline{\hspace{1cm}}, Pl.: \underline{\hspace{1cm}})
    \item der Mensch (Dativ Sg.: \underline{\hspace{1cm}}, Pl.: \underline{\hspace{1cm}})
\end{enumerate}

\chapterend